% Auxiliary input for the parent's prospective dimension-six construction.
% No complete local-geometric assertion is made in this fragment.
\paragraph{Auxiliary Hodge calculation for a quartic containing a line.}
Fix a line $D_0\subset\mathbb P^3_{\mathbb C}$ and let
$B\subset\mathbb P H^0(\mathcal I_{D_0}(4))$ be the nonempty open subset
parametrizing smooth quartics. We choose a \emph{Hodge-generic complex
point} $b\in B$ as part of the data, in the precise sense described below.
Write $W=W_b$, $H=c_1(\mathcal O_W(1))$, and $D=D_0\subset W$.
We can and do choose $b$ so that
\[
 \begin{gathered}
 \operatorname{NS}(W)=\mathbb ZH\oplus\mathbb ZD,
 \qquad
 \begin{pmatrix}H^2&H\cdot D\\H\cdot D&D^2\end{pmatrix}
 =\begin{pmatrix}4&1\\1&-2\end{pmatrix},\\
 T=\operatorname{NS}(W)_{\mathbb Q}^{\perp}
          \subset H^2(W,\mathbb Q),
 \qquad \operatorname{End}_{\mathrm{Hdg}}(T)=\mathbb Q.
 \end{gathered}
\]
Let $C$ be the Fermat plane quartic and put $S=W\times C$.
Then the largest rational sub-Hodge structure of $H^3(S,\mathbb Q)$
whose Hodge types are contained in $\{(2,1),(1,2)\}$ is
\begin{equation}\label{eq:N6-level-one}
       \operatorname{NS}(W)_{\mathbb Q}\otimes H^1(C,\mathbb Q).
\end{equation}
In particular, for every smooth projective complex curve $P$ and every
$z\in\mathrm{CH}^2(W\times C\times P)$, the component of its rational
cycle class in $H^3(S,\mathbb Q)\otimes H^1(P,\mathbb Q)$ belongs to
\[
 \operatorname{NS}(W)_{\mathbb Q}\otimes
 H^1(C,\mathbb Q)\otimes H^1(P,\mathbb Q).
\]

\emph{Choice of the quartic and its Hodge-generic complex point.}\par
Theorem~1.1 and Corollary~1.3(c) of
\citet[p.~2 of the opened author version]{frBrevikNollet2011}
apply to $D_0$ with $d=4$: $\mathcal I_{D_0}(3)$ is globally generated
and $H^0(\mathcal I_{D_0}(2))\ne0$. They give the asserted Picard group
for a very general member of $B$. A smooth quartic is a K3 surface,
and adjunction gives the displayed intersection matrix.

For completeness, we also justify the endomorphism condition rather
than inferring it solely from the Picard rank. Locally mark the second
cohomology. The Local Torelli Theorem and its line-bundle formulation
in \citet[Proposition~6.2.8 and Section~6.2.4, pp.~107--108 of the
opened author version]{frHuybrechtsK3Book} identify the local deformations
preserving $H,D$ with an open subset of the period quadric in
$(\mathbb QH+\mathbb QD)^\perp_{\mathbb C}$. It has dimension $18$.
These are locally realized by quartics containing a line: very ampleness
of $H$ persists; Riemann--Roch, $D^2=-2$, and $H\cdot D=1$ make $D$
effective, and an effective divisor of hyperplane degree one is a line.
A projective coordinate change puts that line back at $D_0$.

For every nonscalar $a\in\operatorname{End}_{\mathbb Q}(T)$, the
condition that $a$ is a Hodge endomorphism implies that the period line
is an eigenline of $a$. The eigenlines form a finite union of proper
linear subspaces. They cannot contain the nondegenerate irreducible
period quadric, which spans the ambient projective space. Thus each
such Hodge-endomorphism locus is proper. There are countably many
rational endomorphisms. Applied to the polarized weight-zero variation
$\operatorname{End}(T)$, the Hodge-locus theorem of
\citet[Theorem~1.1 and Corollary~1.3, p.~2]{frCDK1994}
shows that these loci are algebraic (with their monodromy translates).
Avoiding their countable union, together with the extra Picard loci,
produces the stipulated complex point. This is the meaning of the
Hodge-generic choice. We do not assert that every embedding of a field
with coefficientwise algebraically independent quartic parameters has
this property.

\emph{The product Hodge group and its representations.}
By \citet[Theorem~2.2.1 and Example~2.2.3, p.~211]{frZarhin1983},
\[
                    \operatorname{Hg}(T)=\operatorname{SO}(T).
\]
The exact original theorem and its proof in Section~2.5, pp.~212--213,
have been inspected. By \citet[Lemma~2, p.~106, and Example~3,
pp.~106--107]{frBeauville2014}, the Jacobian of $C$ is isogenous to
$E^3$, where $E=\mathbb C/\mathbb Z[i]$. Hence, for $V=H^1(E,\mathbb Q)$,
we have $H^1(C,\mathbb Q)\simeq V^{\oplus3}$ and
$\operatorname{Hg}(V)=\operatorname{Res}^{1}_{\mathbb Q(i)/\mathbb Q}
\mathbb G_m=:U$.

Let $G=\operatorname{Hg}(T\oplus V)$. Its two projections are surjective,
and $G$ is connected. The image of $G^{\mathrm{der}}$ on
$\operatorname{SO}(T)$ is the full group, while its image on the torus
$U$ is trivial. Therefore $\operatorname{SO}(T)\times1\subset G$;
surjectivity on $U$ then gives
\[
                   G=\operatorname{SO}(T)\times U.
\]
This argument uses only that the special orthogonal group in rank $20$
is connected and semisimple; no independence assumption on periods
of $W$ and $E$ is being inserted.

Over $\mathbb Q(i)$, the representation $T\otimes V$ is the direct sum
of the two twists of the standard $\operatorname{SO}(T)$ representation
by the two opposite nontrivial characters of $U$. Both summands are
absolutely irreducible, they are nonisomorphic, and complex conjugation
exchanges them. A rational invariant subspace therefore contains both
or neither. Thus $T\otimes V$ is a simple rational Hodge structure.
Moreover
\[
                 (T\otimes V)^{3,0}=T^{2,0}\otimes V^{1,0}\ne0.
\]
It consequently has level three. Semisimplicity of polarizable rational
Hodge structures implies that every nonzero sub-Hodge structure of
$(T\otimes V)^{\oplus3}$ has a simple constituent isomorphic to
$T\otimes V$; in particular none has level at most one.

Since $H^1(W)=H^3(W)=0$, Kunneth gives
\[
 H^3(S,\mathbb Q)=
 \bigl(\operatorname{NS}(W)_{\mathbb Q}\otimes H^1(C,\mathbb Q)\bigr)
 \oplus\bigl(T\otimes H^1(C,\mathbb Q)\bigr).
\]
The first summand has exactly the two indicated Hodge types and the
second has no nonzero sub-Hodge structure with those types. This proves
\eqref{eq:N6-level-one}.

Finally, the indicated Kunneth component of the class of $z$ is a
rational Hodge class of type $(2,2)$. Poincare duality on $P$ turns it
into a morphism of Hodge structures
$H^1(P,\mathbb Q)(-1)\to H^3(S,\mathbb Q)$. The source has types
$(2,1),(1,2)$, so its image lies in \eqref{eq:N6-level-one}, proving
the claimed consequence for every such cycle and parameter curve.
